\chapter{Conclusion and Recommendation}

\section{Summary of the Study}
This project set out to design and implement a robotic system capable of recognizing human intention through hand gestures and body movement, and responding autonomously without requiring cloud connectivity or GPU-class hardware. The system was built on a Yahboom MicroROS-Pi5 platform, combining a MediaPipe HandLandmarker-based gesture pipeline, a Multilayer Perceptron (MLP) gesture classifier trained on 19 engineered geometric hand features, a Long Short-Term Memory (LSTM) network for short-horizon movement prediction, and a spatial-zone-based multi-person tracking approach with gesture-confirmation locking. Autonomous navigation infrastructure (SLAM Toolbox, AMCL, Nav2) was integrated into the same pipeline to support movement between human encounters. The physical prototype demonstrated that multi-modal intention recognition can execute reliably on a low-cost, low-power edge computing architecture.

Across the course of the project, six specific objectives were pursued and successfully realized:
\begin{enumerate}
    \item A custom, balanced dataset of 6,000 gesture samples across six discrete commands (Stop, Follow, Go, Left, Right, Back) was collected directly through the onboard monocular camera, eliminating sensor domain shift.
    \item A lightweight Multilayer Perceptron (MLP) gesture classifier was developed and trained on 19 scale-invariant geometric hand features, achieving high classification fidelity.
    \item An LSTM sequential motion prediction model was developed and evaluated to anticipate an operator's short-horizon trajectory from continuous body pose landmarks.
    \item A ROS2-based multi-node software architecture was implemented, integrating vision perception with a centralized decision-making state machine (\texttt{brain\_node}) for deterministic robot control.
    \item A spatial-zone filtering mechanism coupled with a rolling majority-vote confirmation buffer was engineered to maintain active lock on a primary operator and reject background bystander interference.
    \item Simultaneous Localization and Mapping (SLAM) was configured and physically validated on the mobile platform, and comprehensive physical testing was conducted to evaluate classification accuracy, motor response, and hardware throughput.
\end{enumerate}

\section{Conclusion}
The experimental results demonstrate that explicit, gesture-based human-robot command interfaces can be successfully combined with implicit, motion-based intention inference on low-cost edge hardware without cloud dependency. The gesture recognition subsystem achieved a test accuracy of 99.38\% on a balanced, 1,000-sample-per-class dataset of six gesture classes (Stop, Follow, Go, Left, Right, Back). Physical trials across three participants performing each gesture ten times returned per-gesture accuracy ranging from 83\% to 100\%, with five of the six gestures staying within a narrow 7--12 percentage-point band. This confirms that the engineered geometric features successfully eliminate perspective distortion across operational distances.

The movement prediction subsystem, built on an LSTM trained on 152 motion sequences, achieved an Average Displacement Error of 25.8~px and a Final Displacement Error of 32.4~px on held-out test sequences. All five movement classes (Approaching, Moving Away, Moving Left, Moving Right, Stationary) were successfully and consistently recognized during functional testing. The model's lateral position predictions were streamed directly through the \texttt{Detection} message, providing anticipatory context for downstream decision-making.

Physical deployment surfaced and resolved several real engineering constraints not apparent from simulation alone. These included memory-pressure crashes on the original 2GB Raspberry Pi 5 configuration (resolved by upgrading to 8GB), a \texttt{ROS\_DOMAIN\_ID} mismatch silently preventing inter-node communication, and a brittle rule-based gesture classifier that motivated the shift to a trained MLP. SLAM-based occupancy grid mapping was successfully validated on physical hardware, producing a coherent metric map of a real indoor environment following the diagnosis and correction of transform race conditions and yaw fusion drift. Autonomous path planning through the resulting map via Nav2 was architecturally integrated into the software stack, with full physical navigation validation remaining as an area for ongoing investigation.

\section{Limitations of the Study}
Several technical limitations should be considered when interpreting the results of this study:

\begin{enumerate}
    \item \textbf{Participant Sample Size:} Gesture recognition physical trials were conducted with three participants, each performing every gesture ten times (yielding 30 trials per gesture and 180 physical trials in total). While this sample size proved sufficient to demonstrate functional viability, it is too small a cohort to establish comprehensive statistical generalization across wider demographic variations in hand anatomy, skin tones, and motor proficiencies.
    \item \textbf{Absence of Structured Per-Trial Records:} The gesture and movement testing results reported in Chapter 4 were compiled from summary-level evaluation logs. Individual per-trial records were not logged in a structured database at the time those initial trials were conducted. A dedicated logging node (\texttt{test\_logger\_node.py}) was developed during this project specifically to address this gap for future testing sessions, but it was not yet in active use during initial data collection.
    \item \textbf{Autonomous Navigation Not Independently Validated:} While SLAM-based mapping was validated on physical hardware with drift-related issues diagnosed and corrected, autonomous waypoint path-planning through the resulting map using Nav2 had not been demonstrated end-to-end on hardware at the time of writing. A configuration audit of the navigation stack identified and corrected several issues preventing successful bringup; however, complete closed-loop navigation—including dynamic obstacle avoidance and recovery behaviors in populated environments—remains to be validated in future physical trials.
    \item \textbf{Map Reusability and Loop-Closure Tuning:} A systematic diagnosis resolved two root causes of map quality degradation: a transform race condition causing high scan-drop rates, and disabled yaw fusion causing rotational drift. However, two limitations remain: the occupancy grid map is currently saved only as a flattened image/YAML export without serialized SLAM solver states (meaning mapping sessions cannot be resumed or incrementally updated), and loop-closure parameters have not been tuned beyond default settings.
    \item \textbf{Facial Recognition Not Deployed in Runtime Container:} A facial recognition prototype utilizing InsightFace (ArcFace embeddings) was evaluated offline, successfully distinguishing between enrolled operator identities with confidence scores of 0.54--0.57 against an unknown baseline of 0.09. A ROS2-integrated node (\texttt{face\_id\_node}) was written to productionize this capability, but it has not yet been registered as a buildable package entry point, and the heavy InsightFace dependency is not installed in the robot's runtime container. Consequently, facial identification remains unintegrated into the live system.
    \item \textbf{Person-Following Control Stability:} The Follow command's visual servoing behavior steers the entire robot chassis using a proportional controller based on horizontal tracking error. This simple proportional control is susceptible to overshoot and minor heading oscillations when an operator moves rapidly. Furthermore, the camera gimbal is centered once at system startup and does not track independently of chassis locomotion.
    \item \textbf{Hardware-Specific Environmental Validation:} All testing was conducted on a single physical unit under indoor lighting and environmental conditions typical of the test setting. Performance under significantly different ambient lighting (such as bright outdoor sunlight), non-flat floor surfaces, or complex acoustic environments was not evaluated.
\end{enumerate}

\section{Recommendations for Future Work}
Based on the outcomes and limitations identified during this study, the following recommendations are made for future extensions of this research:

\begin{enumerate}
    \item \textbf{Expand Physical Trial Scale and Rigor:} Future testing should employ the dedicated \texttt{test\_logger\_node.py} instrumentation to capture structured, per-trial telemetry across a larger and more diverse participant pool. This will enable statistically robust evaluation across diverse hand shapes, clothing, and ambient lighting conditions.
    \item \textbf{Validate and Tune Autonomous Navigation:} Nav2-based autonomous path-planning should be independently tested and validated end-to-end on the physical robot. This should include verifying local planner trajectory generation, dynamic obstacle avoidance, and costmap recovery behaviors in crowded hallways.
    \item \textbf{Improve Map Persistence and Quality:} Future work should implement serialized SLAM state saving within SLAM Toolbox, allowing mapping sessions to be resumed, merged, and incrementally refined over time. Additionally, loop-closure parameters should be tuned using deliberate loop-closure driving trajectories to enhance geometric consistency in large environments.
    \item \textbf{Deploy and Integrate Facial Recognition:} The offline facial recognition prototype should be containerized and integrated into the live perception pipeline. Fusing facial identity embeddings with spatial-zone filtering will provide persistent operator authentication, eliminating identity ambiguity when multiple individuals enter the operational zone.
    \item \textbf{Upgrade Person-Following Control to Closed-Loop PID:} Replacing the proportional-only visual servoing controller with a tuned Proportional-Integral-Derivative (PID) controller (combined with Kalman-filtered centroid tracking) will reduce tracking overshoot and damp heading oscillations. Independently, enabling active 2-DOF gimbal tracking will allow the camera to center the human operator without requiring continuous chassis rotation.
    \item \textbf{Broaden Environmental and Sensor Validation:} Future research should evaluate the system across varied lighting conditions, surface textures, and physical layouts. Integrating complementary sensing modalities—such as ultrasonic proximity sensors or low-power solid-state depth cameras—would provide redundancy and enhance near-field obstacle detection.
\end{enumerate}
